\section{Self-Attention 与 Transformer}

\subsection{Self-Attention 的四个步骤}

Self-Attention（自注意力）是 Attention 的一种特殊形式：\textbf{Query、Key、Value 都来自同一个输入序列}——即序列对自身的不同位置进行注意力计算。

\begin{figure}[H]
\centering
\includegraphics[width=0.85\textwidth]{slides-images/slide-70.jpg}
\caption{Self-Attention 的四个步骤\protect\footnotemark}
\end{figure}
\footnotetext{Slides 第70页。}

Self-Attention 的完整计算过程分为四步：

\textbf{Step 1: 编码位置信息（Positional Encoding）}

由于我们不再逐步处理序列，而是一次性输入所有元素，我们需要显式地编码位置信息：

$$\tilde{x}_i = x_i + p_i$$

其中 $p_i$ 是位置编码向量，它为每个位置提供唯一的标识。

\begin{importantbox}{Positional Encoding 的必要性}
没有位置编码，Self-Attention 将无法区分 ``the cat sat on the mat'' 和 ``the mat sat on the cat''——因为 Attention 操作本身对输入的顺序是\textbf{置换不变}（Permutation Invariant）的。位置编码打破了这种对称性，让模型能够感知序列中的顺序关系。
\end{importantbox}

\textbf{Step 2: 提取 Query、Key、Value}

对位置编码后的输入，通过三个不同的线性变换层分别生成 Q、K、V 矩阵：

$$Q = \tilde{X} \cdot W_Q, \quad K = \tilde{X} \cdot W_K, \quad V = \tilde{X} \cdot W_V$$

\begin{itemize}
    \item $W_Q, W_K, W_V$ 是三个可学习的权重矩阵
    \item 同一个输入 $\tilde{X}$ 经过不同的线性变换，捕捉不同方面的信息
\end{itemize}

\textbf{Step 3: 计算 Attention 权重}

使用\textbf{缩放点积}（Scaled Dot-Product）计算 Query 和 Key 之间的相似度：

$$\text{Attention Score} = \text{softmax}\left(\frac{Q \cdot K^T}{\sqrt{d_k}}\right)$$

\begin{itemize}
    \item $Q \cdot K^T$：Query 与 Key 的点积，衡量相似度
    \item $\sqrt{d_k}$：缩放因子（$d_k$ 为 Key 的维度），防止点积值过大导致 softmax 梯度消失
    \item $\text{softmax}$：将相似度归一化为概率分布
\end{itemize}

\begin{figure}[H]
\centering
\includegraphics[width=0.65\textwidth]{slides-images/slide-73.jpg}
\caption{使用缩放点积计算相似度\protect\footnotemark}
\end{figure}
\footnotetext{Slides 第73页。}

\textbf{Step 4: 提取高注意力特征}

用 Attention 权重对 Value 进行加权求和，得到最终输出：

$$\text{Output} = \text{softmax}\left(\frac{Q \cdot K^T}{\sqrt{d_k}}\right) \cdot V$$

\begin{importantbox}{Self-Attention 的完整公式}
$$\text{Attention}(Q, K, V) = \text{softmax}\left(\frac{Q K^T}{\sqrt{d_k}}\right) V$$

这个公式被称为 Scaled Dot-Product Attention，是 Transformer 架构的核心计算单元。
\end{importantbox}

\subsection{Attention 权重的可解释性}

Attention 权重矩阵提供了一种直观的可解释性：矩阵中的每个元素 $a_{ij}$ 表示位置 $i$ 对位置 $j$ 的``关注程度''。

\begin{figure}[H]
\centering
\includegraphics[width=0.75\textwidth]{slides-images/slide-75.jpg}
\caption{Attention 权重矩阵的可视化：颜色越深表示关注度越高\protect\footnotemark}
\end{figure}
\footnotetext{Slides 第75页。}

例如，在句子 ``He tossed the tennis ball to serve'' 中，``ball'' 这个词可能对 ``tennis'' 和 ``tossed'' 有较高的注意力权重，因为这些词对于理解 ``ball'' 的语义最为关键。

\subsection{从 Attention Head 到 Transformer}

\begin{figure}[H]
\centering
\includegraphics[width=0.65\textwidth]{slides-images/slide-78.jpg}
\caption{Self-Attention 是 Transformer 架构的基础构建模块\protect\footnotemark}
\end{figure}
\footnotetext{Slides 第78页。}

一个完整的 Self-Attention 计算（从位置编码到 Q/K/V 提取再到加权输出）构成一个 \textbf{Attention Head}。

\subsection{Multi-Head Attention}

实际的 Transformer 不只使用一个 Attention Head，而是使用\textbf{多个}并行的 Attention Head——这就是 \textbf{Multi-Head Attention}：

$$\text{MultiHead}(Q, K, V) = \text{Concat}(\text{head}_1, \ldots, \text{head}_h) W^O$$

其中每个 head 使用不同的 $W_Q^i, W_K^i, W_V^i$：

$$\text{head}_i = \text{Attention}(Q W_Q^i, K W_K^i, V W_V^i)$$

\begin{importantbox}{Multi-Head Attention 的优势}
多头注意力让模型能够\textbf{同时从不同的角度}关注输入的不同方面。例如，在处理一个句子时：
\begin{itemize}
    \item 一个 Head 可能专注于捕捉语法结构关系
    \item 另一个 Head 可能专注于语义相关性
    \item 第三个 Head 可能关注指代消解（谁指代谁）
\end{itemize}
多个 Head 的输出被拼接并通过线性变换融合，使模型能够综合利用多种关注模式。
\end{importantbox}

\subsection{本章小结}

Self-Attention 通过四步操作（位置编码、Q/K/V 提取、相似度计算、加权求和）实现了对序列内部依赖关系的直接建模。它不需要逐步处理序列，可以并行计算，且能直接捕捉任意距离的依赖关系。Multi-Head Attention 进一步增强了模型从多个角度理解输入的能力。这些机制构成了 Transformer 架构的核心。

\newpage
